% Test fixture for audit_tex.py. Every problem below is planted on purpose;
% tests/test_audit_tex.py lists what each check must find.
\documentclass[aps,prl,reprint]{revtex4-2}
\usepackage{graphicx}
\usepackage{xcolor}
\begin{document}

\title{Can strain stabilize skyrmions in YIG?}
\author{A. Author}
\affiliation{XYZ Laboratory, Springfield}

\begin{abstract}
Magnetic skyrmions are candidates for dense memory.
Here we predict that strain stabilizes skyrmions in thin films, widening the stability window by a factor of 2.6 at $T = 300$~K.
The DMI sets the skyrmion size, and the effect reaches 45\% of the bulk value.
\end{abstract}

\maketitle

Spin-orbit torque (SOT) drives skyrmions efficiently~\cite{b,a}.
SOT also moves domain walls.
We integrate the LLG equation.
Figure~\ref{fig:two} shows the stability window, and Fig.~\ref{fig:one} shows the geometry.
Is this effect intrinsic?
The energy density is
\begin{equation}
E = A (\nabla \mathbf{m})^2 ,
\label{eq:energy}
\end{equation}
where $A$ is the exchange stiffness.

\begin{figure}
\includegraphics{f1}
\caption{Geometry. (a) Film; (b) strain profile.}
\label{fig:one}
\end{figure}

\begin{figure}
\includegraphics{f2}
\caption{Is the window wider?}
\label{fig:two}
\end{figure}

\begin{figure}
\includegraphics{f3}
\caption{Never cited.}
\label{fig:three}
\end{figure}

The stability window grows by a factor of 2.6 at 300~K.
\textcolor{red}{\textbf{[TODO: data for the thick-film case]}}
See also Ref.~\cite{c} and Fig.~\ref{fig:missing}.
% TODO: ask about the substrate

\begin{thebibliography}{9}
\bibitem{a} A. One, Phys. Rev. Lett. \textbf{1}, 1 (2020).
\bibitem{b} B. Two, Phys. Rev. B \textbf{2}, 2 (2021).
\bibitem{d} D. Four, Nat. Phys. \textbf{4}, 4 (2022).
\end{thebibliography}

\end{document}
